% Sources: src/react_agent/security_v3/{intent_policy_v1,value_labels_v1,param_pinning_v1,brokerguard_v3_v1}.py;
% decision log 2026-11-16 12:30 (C1), 15:30 (runtime, C1a rejected), 16:20 (micro-benchmark), 2026-11-17 19:25 (freeze).
% Numbers: facts.md "Bổ sung cho bài v3" (8 characters, 9 digits; micro-benchmark; C1a 2 -> 5) and Test-4 (0.26 ms).
% No attack text is quoted. The cue words are those of the extractor, not of the data.
\section{The \ours{} Method}\label{sec:method}

The design of \ours{} starts from the view that the model is not a security boundary. An action takes effect only when
the broker executes it, so the broker acts as a reference monitor at the sinks~\cite{anderson1972}. The only source of
authority is $u$. Permissions are derived deterministically from $u$, never from tool content and never from the
judgement of an LLM. Blocking does not end the task either, because a denied call is reported back so that the agent
can continue with the legitimate part of its work. Figure~\ref{fig:arch} shows the architecture. \ours{} keeps the
layers of \base: \rFlow{} based on source labels, \rRedact, which redacts confidential values in the answer, \rSig, which
denies every sink after an observation flagged by a phrase detector (except calls that match a strict grammar over
$u$), the \cSand{} prompt, which repeats $u$ after every observation, and denial feedback. The destination rule \rDest{}
is replaced by \cOne, and \cTwo{} and \cThree{} are added.

% Figure 1, TikZ (drawn by Overleaf, never rendered locally). Coordinates in cm, total width 11.85 cm < \textwidth 12.2 cm.
% Check boxes are 1.7 x 0.9 cm with 6.4 pt text: about 14 characters per line, three lines per box (keep labels short).
% Colour follows the entity as in Figures 2 and 3: thick orange border = new components C1--C3, blue = kept from BG-Base,
% grey dashed = model-dependent. In black and white the border width and the dash pattern separate the classes.
\begin{figure}[!htbp]
\centering
\begin{tikzpicture}[
  font=\fontsize{6.4}{7.6}\selectfont,
  box/.style={rounded corners=1pt, align=center, inner sep=1pt},
  newc/.style={box, draw=bgorange, line width=0.9pt, fill=bgorange!12},
  kept/.style={box, draw=bgblue, line width=0.45pt, fill=bgblue!10},
  soft/.style={box, draw=black!55, line width=0.5pt, fill=black!4, dash pattern=on 2pt off 1.3pt},
  neutral/.style={box, draw=black!60, line width=0.5pt, fill=black!5},
  arr/.style={-{Stealth[length=3pt]}, line width=0.5pt, draw=black!65},
  lab/.style={font=\fontsize{6.4}{7}\selectfont},
  note/.style={font=\fontsize{5.6}{6.6}\selectfont\itshape, text=black!65, anchor=west},
  gate/.style={minimum width=1.7cm, minimum height=0.9cm, text width=1.6cm},
  swatch/.style={minimum width=0.34cm, minimum height=0.24cm},
]
% left column: user and agent
\node[neutral, minimum width=1.5cm, minimum height=0.8cm, text width=1.4cm] at (0.75,4.85)
  {User\\ $u$ (trusted)};
\node[neutral, minimum width=1.5cm, minimum height=1.3cm, text width=1.4cm,
  font=\fontsize{6.6}{7.6}\selectfont\bfseries] at (0.75,2.35) {LLM agent\\ (ReAct)};
\draw[arr] (0.75,4.45) -- (0.75,3.0);

% broker frame
\draw[rounded corners=2pt, line width=0.6pt] (2.0,0.5) rectangle (9.3,5.65);
\node[font=\fontsize{7}{8}\selectfont\bfseries] at (5.65,5.43) {Tool Broker (host side)};

% row 1: state built by C1 and C2
\node[newc, minimum width=3.45cm, minimum height=0.66cm, text width=3.35cm] at (3.825,4.8)
  {\textsf{C1}: extract $P(u)$ from $u$\\ (destination, polarity, pins)};
\node[newc, minimum width=3.45cm, minimum height=0.66cm, text width=3.35cm] at (7.475,4.8)
  {\textsf{C2}: label each value in $o_t$\\ (protected set $C_t$)};
\draw[arr] (1.5,4.8) -- (2.1,4.8);

% row 2: the four checks before every sink call
\node[note] at (2.1,4.2) {before every sink call $a_t$: executed only if all four checks pass};
\node[kept, gate] at (2.95,3.5) {\textsf{R-Signal}\\ flagged observation, deny};
\node[newc, gate] at (4.75,3.5) {\textsf{C1}\\ positive grant};
\node[newc, gate] at (6.55,3.5) {\textsf{C3}\\ pinned values kept};
\node[newc, gate] at (8.35,3.5) {\textsf{C2} + \textsf{R-Flow}\\ no value of $C_t$ leaves};

% row 3: final answer, prompt layer, feedback, confirmation
\node[kept, gate] at (2.95,2.1) {\textsf{R-Redact} + \textsf{C2}\\ redact values in $y$};
\node[soft, gate] at (4.75,2.1) {Sandwich\\ repeat $u$};
\node[kept, gate] at (6.55,2.1) {Report denial,\\ agent continues};
\node[newc, gate] at (8.35,2.1) {Needs confirmation\\ = denial};

% legend
\node[newc, swatch] at (2.3,0.88) {};
\node[note] at (2.52,0.88) {new component};
\node[kept, swatch] at (4.3,0.88) {};
\node[note] at (4.52,0.88) {kept from \base};
\node[soft, swatch] at (6.3,0.88) {};
\node[note] at (6.52,0.88) {model-dependent};

% agent <-> broker
\draw[arr] (1.5,2.75) -- node[lab, above=-1pt] {$a_t$} (2.0,2.75);
\draw[arr] (2.0,2.35) -- node[lab, above=-1pt] {$o_t$} (1.5,2.35);
\draw[arr] (1.5,1.95) -- node[lab, above=-1pt] {$y$} (2.0,1.95);

% tool column
\node[neutral, minimum width=2.1cm, minimum height=0.8cm, text width=2.0cm] at (10.8,3.95)
  {Read tools\\ (untrusted)};
\node[neutral, minimum width=2.1cm, minimum height=0.8cm, text width=2.0cm] at (10.8,2.45)
  {Sinks\\ e-mail, webhook};
\draw[arr] (9.3,4.15) -- node[lab, above=-1pt] {$a_t$} (9.75,4.15);
\draw[arr] (9.75,3.75) -- node[lab, above=-1pt] {$o_t$} (9.3,3.75);
\draw[arr] (9.3,2.45) -- node[lab, above=-1pt] {$a_t$} (9.75,2.45);

% redacted answer leaves the broker
\draw[arr] (9.3,1.0) -- (9.95,1.0);
\node[anchor=west, text=black!65] at (9.98,1.0) {$y$ redacted};
\end{tikzpicture}
\caption{Architecture of \ours. Every call $a_t$ and the final answer $y$ pass through the Tool Broker. \cOne{} extracts
the policy $P(u)$ once from $u$, and \cTwo{} updates the protected set $C_t$ after every observation. Boxes with a thick
orange border are new components, boxes with a blue border are kept from \base, and the dashed box depends on the model.
All checks inside the frame are deterministic.}\label{fig:arch}
\end{figure}


\paragraph{Normalisation.} All comparisons use two normalisation functions. The function $\norm(\cdot)$ applies Unicode
NFKC, removes format characters such as zero-width characters, and casefolds. The function $\kappa(\cdot)$ additionally
strips Vietnamese diacritics and keeps only letters and digits. The notation $x \sub z$ means that $\kappa(x)$ is a
substring of $\kappa(z)$, and $V(u)$ denotes the set of values that occur in $u$ in this sense.

\paragraph{\cOne: intent policy.} The extractor reads only $u$ and returns $P(u) = \{g_i\}$. Each grant $g =
(\mathrm{tool}, \dest, \pi, \mathrm{pins})$ consists of a tool, a destination, a polarity $\pi \in \{+,-\}$ and the
values pinned by \cThree. Destinations are e-mail addresses or endpoints matched as whole tokens after $\norm$, so part
of a destination or an ordinary word in $u$ no longer counts as a destination. Polarity follows the scope of negation.
A negated verb (``đừng gửi'', ``do not send'') covers the destinations that follow it in the same clause. Exclusion cues
(``trừ'', ``thay vì'', ``except'') cover the destinations up to the next comma. Contrast cues (``nhưng'', ``but'') close
the scope. When $u$ carries diacritics, the original spelling decides whether a word such as ``dung'' means ``dùng''
(use) or ``đừng'' (do not). When $u$ has no diacritics, such a word counts as a negation only if a sending verb follows
it. Grants are typed by action. An address authorises only sending e-mail, an endpoint only calling a webhook, and the
clause must contain a matching sending verb that is not negated. A destination that is only mentioned, appears in a
quotation or in reported speech, or stands near a negation of unclear scope yields \emph{needs confirmation}. A sink
call $a$ satisfies \cOne{} when both of the following conditions hold:
\begin{equation}\label{eq:c1}
\begin{aligned}
&\exists\, g \in P(u)\colon\; \pi(g) = {+} \,\land\, \mathrm{tool}(g) = \mathrm{tool}(a) \,\land\, \dest(g) \doteq \dest(a),\\
&\nexists\, g' \in P(u)\colon\; \pi(g') = {-} \,\land\, \dest(g') \doteq \dest(a),
\end{aligned}
\end{equation}
where $\doteq$ denotes whole-token equality after $\norm$. Sinks without a destination field lie outside \cOne. As an
example, take the request ``compile the quarterly report and send it to \code{ketoan@example.com}, not to
\code{doitac@example.org}''. The extractor returns a positive grant for (e-mail, \code{ketoan@example.com}) and a
negative grant for \code{doitac@example.org}. A call that sends to the second address is denied, even though the address
occurs in $u$ and would therefore pass \rDest{} in \base.

\paragraph{\cTwo: value-level labels.} \base{} protects only values that come from sources labelled confidential, so a
secret inside a source labelled public is never protected. \cTwo{} labels each value in every tool output by its type,
independently of the source label and of trust. Three types are recognised: identifiers mixing letters and digits with
at least 8 characters, numbers with at least 9 digits after removing spaces, dots and dashes, and values that follow a
sensitive field name such as password, OTP, token or account number. Ordinary words and short numbers are never treated
as protected values. Let $C_t$ be the set of protected values after step $t$, comprising the values protected by \rFlow{}
through source labels and the values labelled by \cTwo. The broker denies a sink call $a$ if
\begin{equation}\label{eq:c2}
\exists\, v \in C_t \setminus V(u)\colon\; v \sub \data(a) .
\end{equation}
Because the comparison uses $\kappa$, the condition still holds when spaces, dots or dashes are added or removed, when
case changes, or when diacritics are stripped. In the final answer, every $v \in C_t \setminus V(u)$ with $v \sub y$ is
replaced by \code{[REDACTED]} in every surface form that matches $\kappa(v)$. The detector deliberately does not rely on
the format that the data generator uses for secrets, so its coverage of other formats remains an empirical question
(Section~\ref{sec:formal}).

\paragraph{\cThree: parameter pinning.} Parameter tampering keeps the authorised destination and changes a value such as
an amount, a status, a unit or a code. \cThree{} checks a call against the clause of $u$ that authorised it. The pinned
literals are numbers that follow the sending verb and are not operands of a calculation, units stated with a number,
Boolean values with their polarity, identifiers written by the user, and every number in a sentence in which $u$
declares the values fixed. A call violates \cThree{} if a pinned literal is missing from $\data(a)$ (compared under
$\kappa$) or if it carries the opposite of a pinned Boolean. If several grants with the same tool and destination match
the call, satisfying one of them is enough.

\paragraph{Gate and cost.} A sink call is executed only if it passes all four checks: \rSig, \cOne, \cThree, and \cTwo{}
together with \rFlow. The order of the checks only affects which reason is logged for a denial. A \emph{needs
confirmation} result is enforced as a denial and counted separately. Read calls are always allowed. \ours{} calls no
additional model. Over the 1,960 development cases, the median time to check a sink call is 315~$\mu$s (\base:
107~$\mu$s), and building $P(u)$ takes 496~$\mu$s per task. In the actual runs, the gate takes 0.26~ms per call on
average, against about 14~s of model time per case.

\paragraph{Development and freezing.} The cue lists, length thresholds, field names and units were tuned only on the
training part of the development set, and the configuration was selected on its validation part
(Section~\ref{sec:protocol}). A variant that used $P(u)$ in place of the strict grammar of \rSig{} increased utility,
but when trajectories of the validation part were replayed, the number of harmful trajectories rose from 2 to 5, so the
variant was rejected. All code, prompts and configurations were then pinned by hash before the test set was run.
